% ============================================================
%  EXAMINER STRATEGY GUIDE
%  PP-VAE MPhil Dissertation — USN 4730
%  For Overleaf: compile as standalone or include in main.tex appendix
% ============================================================

\documentclass[12pt, a4paper]{article}
\usepackage[margin=2.5cm]{geometry}
\usepackage{libertine}
\usepackage[T1]{fontenc}
\usepackage{booktabs}
\usepackage{longtable}
\usepackage{array}
\usepackage{xcolor}
\usepackage{hyperref}
\usepackage{enumitem}
\usepackage{mdframed}

\definecolor{darkblue}{RGB}{0,51,102}
\definecolor{warnyellow}{RGB}{200,150,0}
\definecolor{okgreen}{RGB}{0,120,60}
\definecolor{critred}{RGB}{180,30,30}

\title{\textbf{\color{darkblue}Distinction-Level Strategy Guide}\\
\large Pathology-Preserving VAE Dissertation — USN 4730\\
Cambridge MPhil Population Health Sciences}
\date{Confidential working document}

\begin{document}
\maketitle
\tableofcontents
\newpage

% ============================================================
\section{Potential Examiner Criticisms}
% ============================================================

\noindent The following criticisms represent every substantive objection
a Cambridge examiner is likely to raise. Each is ranked by probability
of occurrence (H/M/L) and severity if unaddressed (H/M/L).

\begin{longtable}{p{0.5\linewidth} p{0.08\linewidth} p{0.08\linewidth}}
\toprule
\textbf{Criticism} & \textbf{Prob.} & \textbf{Sev.} \\
\midrule
\endhead

\textbf{C1: Resolution triviality.}
``28×28 images are so small that any deep learning model
will reconstruct them adequately. The results are an artefact
of the easy benchmark, not evidence of a methodological advance.''
& H & H \\

\addlinespace
\textbf{C2: Baseline VAE outperforms PP-VAE on PSNR.}
``If the proposed model achieves lower PSNR than the simpler
baseline, the extra complexity is unjustified.''
& H & H \\

\addlinespace
\textbf{C3: No radiologist validation.}
``PSNR and SSIM do not measure diagnostic utility. Without a
reader study, clinical claims are unsupported.''
& H & M \\

\addlinespace
\textbf{C4: Loss weight selection is arbitrary.}
``The lambda values appear to have been tuned post-hoc to
maximise performance, biasing the ablation results.''
& H & M \\

\addlinespace
\textbf{C5: Perceptual loss domain mismatch.}
``VGG-16 features trained on ImageNet are inappropriate
for greyscale 28×28 medical images.''
& M & M \\

\addlinespace
\textbf{C6: No comparison to diffusion models.}
``The current SOTA is DDPM-based. Without that comparison
the work's positioning is incomplete.''
& M & M \\

\addlinespace
\textbf{C7: Single-dataset evaluation.}
``Results on PneumoniaMNIST cannot be generalised to any
clinical imaging setting.''
& H & M \\

\addlinespace
\textbf{C8: Statistical testing on 624 images is underpowered
for subgroup analyses.}
& M & L \\

\addlinespace
\textbf{C9: Novelty is incremental.}
``Combining existing loss terms is engineering, not science.''
& M & H \\

\addlinespace
\textbf{C10: Calibration metrics not validated.}
``ECE is defined for classification; its use for pixel-wise
reconstruction uncertainty is non-standard.''
& M & M \\

\addlinespace
\textbf{C11: Class imbalance not addressed.}
``73\% pneumonia may inflate aggregate metrics on the
majority class.''
& M & L \\

\addlinespace
\textbf{C12: Simulated degradation is not clinically validated.}
``The Poisson-Gaussian model is theoretical; real LDCT
physics include additional artefacts.''
& M & M \\

\addlinespace
\textbf{C13: Reproducibility not demonstrated.}
``Results require independent reproduction before claims
can be accepted.''
& L & M \\

\addlinespace
\textbf{C14: Lack of interpretability beyond t-SNE/UMAP.}
``Visualisations do not constitute mechanistic understanding.''
& L & L \\

\bottomrule
\end{longtable}

% ============================================================
\section{Evidence Required to Neutralise Each Criticism}
% ============================================================

\begin{description}[style=nextline, leftmargin=0pt]

\item[\color{critred}C1 — Resolution triviality]
\textit{Pre-empt in Section 3.2.2 (Dataset Justification):}
(a) Cite sharpXR (Abolade et al. 2025) showing that essential
anatomical features — cardiac silhouette width, lung field
opacification, costophrenic angle blunting — are distinguishable
at 28×28 in a clinical workflow analogue.
(b) Frame explicitly as ``architectural validation'': the goal is
to characterise the relative contribution of loss terms under
controlled conditions, not to deploy.
(c) Show that class-discriminative features (used in downstream
classification) are preserved after reconstruction — if a ResNet
trained on clean images and fine-tuned on reconstructed images
achieves similar AUC, anatomical content is preserved.

\item[\color{critred}C2 — Baseline VAE beats PP-VAE on PSNR]
\textit{Address in Sections 4.3.1 and 5.1.1:}
(a) Explain the PSNR–SSIM trade-off: MSE-optimised models
regress to the conditional mean, yielding artificially high PSNR
but blurred structure. PP-VAE's constraints suppress mean regression,
reducing PSNR while improving SSIM and perceptual fidelity.
(b) Show that SSIM improvement is statistically significant and
clinically more meaningful for radiograph interpretation.
(c) Include a difference map demonstrating that baseline VAE's
higher PSNR comes from blurring pathological texture, not
preserving it.
(d) If a perceptual metric (LPIPS) was computed, show PP-VAE wins.

\item[\color{critred}C3 — No radiologist validation]
\textit{Address in Sections 5.4 and 5.6:}
(a) Acknowledge explicitly as the primary limitation.
(b) Cite precedent: major VAE reconstruction papers (Ehrhardt
\& Wilms 2022; Zhao \& Huang 2025) similarly rely on computational
metrics in their initial publications.
(c) Propose a concrete reader study protocol as Future Work (F2).
(d) Note that MedMNIST classification labels provide a weak proxy:
if reconstructions from PP-VAE support downstream classification
at least as well as baseline VAE, diagnostic content is preserved.

\item[\color{critred}C4 — Arbitrary loss weights]
\textit{Address in Sections 3.6.7 and Appendix B:}
(a) Document the weight selection procedure: grid search or random
search over a pre-defined range on validation set only.
(b) Show sensitivity analysis (Section 4.6.1) demonstrating
that performance is stable across a range of lambda values
— the result is not critically dependent on exact values.
(c) Fix weights before test set evaluation and report that the
test set was accessed once only.

\item[\color{critred}C5 — Perceptual loss domain mismatch]
\textit{Address in Section 3.6.5:}
(a) Cite Johnson et al. (2016) and Dosovitskiy \& Brox (2016)
showing that early VGG layers capture orientation, frequency, and
texture statistics that transfer across domains.
(b) Use only early VGG layers (relu1\_2, relu2\_2) which capture
edge and texture, not semantic content.
(c) If feasible, compare against a medical-imaging-specific feature
extractor (e.g., features from a ResNet trained on ChestX-ray14)
— include as a sensitivity arm or note as future work.

\item[\color{critred}C6 — No diffusion comparison]
\textit{Address in Sections 2.3.3 and 5.4:}
(a) Explicitly justify exclusion in Methods: DDPM inference
requires O(T) steps (T~1000); inference latency is clinically
impractical on CPU-only hardware targeted for LMIC deployment.
(b) Reference compute budget as constraint; note this is also
a limitation acknowledged in sharpXR.
(c) Frame PP-VAE as the ``deployable'' point of the
quality–compute–uncertainty frontier, not the perceptual quality
maximiser.

\item[\color{critred}C7 — Single-dataset evaluation]
\textit{Address in Section 5.7 (Generalisability):}
(a) Acknowledge as a primary limitation.
(b) Argue that PneumoniaMNIST is the appropriate benchmark
for paediatric CXR proof-of-concept (peer-reviewed, widely used).
(c) Describe MIMIC-CXR / CheXpert validation as the immediate
next step.

\item[\color{critred}C8 — Subgroup power]
\textit{Address in Section 3.8:}
Report post-hoc power analysis. If power is low for subgroups,
acknowledge and interpret subgroup results as exploratory.

\item[\color{critred}C9 — Incremental novelty]
\textit{Address in Sections 1.6 and 5.2:}
(a) Novelty claim is the \textit{combination and systematic ablation}
in the paediatric CXR context: no prior work has done this.
(b) The five constraint design with KL annealing is a coherent
methodological package, not ad-hoc combination.
(c) The reproducible benchmark pipeline is itself a contribution
to the field.

\item[\color{critred}C10 — Calibration metric non-standard]
\textit{Address in Sections 3.7.3 and 5.1.4:}
(a) Define the calibration metric precisely: pixel-level uncertainty
vs. reconstruction error bins — analogous to classification ECE.
(b) Cite Kendall \& Gal (2017) and Ovadia et al. (2019) for the
framework of evaluating model uncertainty against empirical error.
(c) Acknowledge the non-standard adaptation and describe it as
an exploratory analysis.

\item[\color{critred}C11 — Class imbalance]
\textit{Address in Section 3.2.1 and Appendix D:}
Report all aggregate metrics separately for Pneumonia and Normal.
Show that improvements are consistent across both classes.

\item[\color{critred}C12 — Simulated degradation]
\textit{Address in Sections 2.2 and 5.6:}
(a) Cite Foi et al. (2008) R²>0.99 validation.
(b) Cite Lee et al. (2018) clinical LDCT validation of the
Poisson-Gaussian model.
(c) Acknowledge the sinogram-domain physics gap as a limitation.

\item[\color{critred}C13 — Reproducibility]
\textit{Address via Appendix F (Checklist) and G (Code):}
Fixed seeds, public data, public code, full hyperparameter table.
Follow Pineau et al. (2021) ML Reproducibility Checklist.

\item[\color{critred}C14 — Interpretability beyond visualisation]
\textit{Address in Section 4.4:}
Supplement t-SNE/UMAP with quantitative metrics (silhouette score,
Davies-Bouldin) and latent traversal analysis demonstrating
continuous, smooth latent structure.

\end{description}

% ============================================================
\section{Missing Analyses That Would Increase Marks}
% ============================================================

\noindent Ranked by expected mark impact (1 = highest):

\begin{enumerate}

\item \textbf{Downstream classification experiment.}
Train a small classifier on (a) clean images, (b) baseline VAE
reconstructions, (c) PP-VAE reconstructions. Compare AUC for
pneumonia detection. If PP-VAE reconstructions enable classification
AUC closer to the clean-image upper bound, this is the strongest
possible argument for pathology preservation. Impact: converts
the dissertation from a ``reconstruction quality'' paper to a
``diagnostic utility'' paper. \textit{Feasible in <1 day of compute.}

\item \textbf{LPIPS (Learned Perceptual Image Patch Similarity) metric.}
pip install lpips. This metric was designed specifically to
correlate with human perceptual judgement. Including it directly
addresses Criticism C2 (PSNR vs. SSIM trade-off) and provides
a peer-reviewed third metric beyond PSNR/SSIM.
\textit{Feasible in <2 hours.}

\item \textbf{Monte Carlo uncertainty visualisation.}
Sample z 30 times for each test image from the PP-VAE encoder
posterior. Compute pixel-wise variance across 30 reconstructions.
Visualise as an uncertainty map. Compare with decoder variance.
Provides evidence of meaningful probabilistic output.
\textit{Feasible in <1 day.}

\item \textbf{KL annealing ablation.}
Explicitly compare: (a) no annealing, (b) linear annealing,
(c) cyclical annealing. If you already ran this, report it.
This directly addresses the ``posterior collapse'' concern and
demonstrates understanding of VAE training dynamics.

\item \textbf{Latent space interpolation figure.}
Take 5 pairs of test images (2 Pneumonia, 2 Normal, 1 cross-class).
Decode 7 equally-spaced interpolated latent codes between each pair.
Show 5×7 grid. A smooth, clinically plausible interpolation is
strong qualitative evidence for a well-structured latent space.
\textit{Feasible in <2 hours.}

\item \textbf{Subgroup stratified results.}
Report PSNR/SSIM separately for Pneumonia vs. Normal across
all ablation arms. Show that PP-VAE improves specifically on
pathological cases. This is the ``pathology-preserving'' claim
made concrete.

\item \textbf{Error histogram analysis.}
Plot the distribution (not just mean) of per-image PSNR and SSIM
for baseline VAE and PP-VAE. Show that PP-VAE reduces the
left tail (badly reconstructed images) even if mean PSNR is lower.
This is a clinically more meaningful metric — worst-case failures
matter most in diagnostic imaging.

\item \textbf{Silhouette score as a function of training epoch.}
Track cluster separation over training for CAE, VAE, and PP-VAE.
Shows that PP-VAE achieves earlier and higher latent separation.
Strong evidence for richer representational learning.

\item \textbf{Computational efficiency table.}
Report: parameters per model, training time per epoch, inference
time per image, memory footprint. Contextualises PP-VAE for
deployment — relevant for examiner interested in clinical translation.

\item \textbf{Failure mode taxonomy.}
Classify failure cases into categories: (a) catastrophic (anatomy
corrupted), (b) moderate (texture wrong), (c) minor (noise
residual). Show counts per category per model. Demonstrates
methodological rigour in error analysis.

\end{enumerate}

% ============================================================
\section{Figures Most Likely to Impress Examiners}
% ============================================================

\noindent Ranked by expected examiner impact:

\begin{enumerate}

\item \textbf{The Six-Panel Reconstruction Figure.}
Columns: degraded input $\mid$ CAE $\mid$ baseline VAE $\mid$ PP-VAE $\mid$ clean ground truth.
Rows: Pneumonia case, Normal case, and the worst-case failure for PP-VAE.
This is the dissertation's centrepiece. Examiners will spend the most
time here. Ensure high visual quality; use \texttt{imshow} with
correct greyscale normalisation.

\item \textbf{The Ablation Bar Chart.}
Side-by-side horizontal bars for PSNR and SSIM, one bar per ablation arm,
colour-coded by which constraints are active. Add significance brackets.
Communicates the systematic nature of the experiment at a glance.

\item \textbf{Latent Space t-SNE Triptych.}
Three panels (CAE $\mid$ VAE $\mid$ PP-VAE) with identical colour coding
(red = Pneumonia, blue = Normal). Shows progression from chaotic to
structured latent space. Quantitative silhouette scores in panel titles.

\item \textbf{Uncertainty Map on a Pathological Case.}
Left: degraded input. Centre: PP-VAE reconstruction. Right: uncertainty
map (decoder variance). High uncertainty should co-localise with the
consolidation region. This is the single most compelling clinical argument
for probabilistic reconstruction.

\item \textbf{KL Annealing and Convergence Figure.}
Two-panel: (a) KL loss trajectory across epochs for all VAE arms,
(b) reconstruction loss (MSE) trajectory. Shows that annealing resolves
posterior collapse and that PP-VAE converges stably.

\item \textbf{Degradation Level Performance Curve.}
PSNR (y) vs. $\eta$ (x) with one line per model. PP-VAE should show
particular advantage at low $\eta$ (extreme dose reduction). This
directly supports the clinical argument.

\item \textbf{Latent Traversal Grid.}
$5 \times 7$ grid of decoded images as latent codes interpolate between
image pairs. Smooth transitions demonstrate a continuous, well-organised
prior. Caption should name the anatomical features that transition smoothly.

\item \textbf{Residual Error Maps.}
Three panels: (a) degraded $-$ clean, (b) baseline VAE $-$ clean,
(c) PP-VAE $-$ clean. Using a diverging colourmap (blue = underestimate,
red = overestimate). Shows that PP-VAE's errors are systematically
smaller in anatomically important regions.

\item \textbf{Downstream Classification AUC Comparison.}
If analysis \#1 above is run: ROC curves for pneumonia detection
on clean / VAE-reconstructed / PP-VAE-reconstructed images.
Demonstrates diagnostic utility directly.

\item \textbf{Four-Phase Pipeline Figure.}
A clean, professionally typeset diagram of the experimental pipeline
(Phase 1 data $\rightarrow$ Phase 2 architecture $\rightarrow$
Phase 3 training $\rightarrow$ Phase 4 evaluation). More polished
than the research protocol version. Use TikZ or draw.io export.

\end{enumerate}

% ============================================================
\section{Narrative Arc: How the Dissertation Tells a Scientific Story}
% ============================================================

\begin{mdframed}[linecolor=darkblue, linewidth=2pt]
\textbf{The Through-Line:}
A child needs a chest X-ray. The radiation that keeps the image sharp
enough to read also increases lifetime cancer risk. We cannot solve this
with hardware alone. The question this dissertation asks is: \textit{given
an image that is already too noisy to read safely, can a probabilistic
deep learning model reconstruct what the clinician needs to see — without
inventing pathology that was never there?}
\end{mdframed}

\bigskip

\noindent\textbf{Chapter 1} opens with the patient. The radiation burden
is quantified in lives (attributable cancers per 10,000 exposures).
The gap between what existing DL reconstruction offers and what it fails
to guarantee — pathological fidelity, calibrated uncertainty — is the
wound in the literature that PP-VAE is designed to heal.

\bigskip

\noindent\textbf{Chapter 2} traces the intellectual lineage. The examiner
should finish reading and think: ``every section of this review was
pointing toward exactly this design.'' The review is not a catalogue;
it is a funnel. Classical denoising established the importance of structure.
Supervised CNNs solved reconstruction but created hallucination risk.
GANs improved perception but destroyed probabilistic interpretability.
Diffusion models achieved the best quality but at undeployable computational
cost. VAEs are the natural home for a deployable, probabilistic,
structurally-aware model. PP-VAE's five constraints are the synthesis
of everything that preceded them.

\bigskip

\noindent\textbf{Chapter 3} is transparent and reproducible. Every design
choice is justified — not just stated. The ablation design is pre-specified.
The statistical analysis plan is declared before results are shown.
The examiner must be able to re-run the experiment from Chapter 3 alone.

\bigskip

\noindent\textbf{Chapter 4} presents a scientific hierarchy. First, validate
the input (degradation pipeline produces clinically realistic noise).
Second, confirm models converge (training curves; no posterior collapse).
Third, report reconstruction (the headline PSNR/SSIM/MAE numbers).
Fourth, explain why it works (latent space; ablation). Fifth, characterise
what we still cannot explain (error analysis; failure cases). The chapter
ends having answered every objective from Chapter 1.

\bigskip

\noindent\textbf{Chapter 5} earns the distinction mark. It does not
simply restate results. It explains \textit{why} PP-VAE's PSNR is lower
than baseline VAE and why this is the expected and correct outcome of
structure-preserving constraints. It situates the work honestly: this is
a proof-of-concept on a benchmark dataset. It names every limitation before
the examiner can raise it. It proposes a future reader study that would
convert the computational finding into clinical evidence.

\bigskip

\noindent\textbf{Chapter 6} closes the loop. The last paragraph returns to
the child in the clinic. It does not overclaim. It states precisely what
is now known that was not known before, and what remains to be shown before
this model could be placed in a radiologist's workflow.

\bigskip

\noindent\textbf{The tone throughout}: confident but not arrogant. The work
is original, systematically conducted, and honestly reported. Every claim
is supported by evidence. Every limitation is acknowledged before
it can be weaponised. The examiner should finish reading thinking:
\textit{this candidate understands not just what they did, but why it
matters and what its boundaries are.} That is the distinction mindset.

\end{document}